\documentclass[12pt,oneside]{book}

% ============================================================
% METADATOS
% ============================================================
\newcommand{\universidad}{Universidad de Buenos Aires (UBA)}
\newcommand{\materia}{Procedimiento Civil}
\newcommand{\titulo}{Notificaciones: Cédulas de Notificación y Procedimiento de Reintentos}
\newcommand{\autor}{Isaac Edgar Camacho Ocampo}
\newcommand{\anio}{2026}

% ============================================================
% PAQUETES
% ============================================================
\usepackage[utf8]{inputenc}
\usepackage[T1]{fontenc}
\usepackage[spanish,es-nodecimaldot]{babel}

\usepackage[
    top=2.5cm,
    bottom=2.5cm,
    left=3cm,
    right=2.5cm,
    headheight=15pt
]{geometry}

\usepackage{amsmath}
\usepackage{amssymb}
\usepackage{mathtools}

\usepackage{graphicx}
\usepackage{float}
\usepackage{caption}
\usepackage{subcaption}

\usepackage{booktabs}
\usepackage{array}
\usepackage{longtable}
\usepackage{tabularx}

\usepackage{enumitem}

\usepackage{xcolor}
\usepackage[most]{tcolorbox}

\usepackage{fancyhdr}
\usepackage{ragged2e}

\usepackage[
    colorlinks=true,
    linkcolor=blue!50!black,
    citecolor=green!40!black,
    urlcolor=blue!60!black,
    pdftitle={\titulo},
    pdfauthor={\autor},
    pdfsubject={Apuntes universitarios},
    bookmarks=true
]{hyperref}

% ============================================================
% CONFIGURACIÓN
% ============================================================
\graphicspath{
    {./img/}
    {./graficos/}
    {./figuras/}
}

\setcounter{tocdepth}{2}

\setlength{\parindent}{0pt}
\setlength{\parskip}{0.5em}

\setlist[itemize]{
    topsep=0.3em,
    itemsep=0.2em
}

\setlist[enumerate]{
    topsep=0.3em,
    itemsep=0.2em
}

\clubpenalty=10000
\widowpenalty=10000

\pagestyle{fancy}
\fancyhf{}
\fancyhead[L]{\nouppercase{\leftmark}}
\fancyhead[R]{\materia}
\fancyfoot[C]{\thepage}
\renewcommand{\headrulewidth}{0.4pt}
\renewcommand{\footrulewidth}{0pt}

\fancypagestyle{plain}{
    \fancyhf{}
    \fancyfoot[C]{\thepage}
    \renewcommand{\headrulewidth}{0pt}
}

% ============================================================
% COMANDOS
% ============================================================
\newcommand{\abs}[1]{\left|#1\right|}
\newcommand{\norm}[1]{\left\lVert#1\right\rVert}
\newcommand{\vect}[1]{\mathbf{#1}}
\newcommand{\deriv}[2]{\frac{d#1}{d#2}}
\newcommand{\pderiv}[2]{\frac{\partial#1}{\partial#2}}

% ============================================================
% ENTORNOS / CAJAS
% ============================================================
\newtcolorbox{definition}[1]{
    enhanced,
    breakable,
    title={Definición: #1},
    fonttitle=\bfseries,
    colback=blue!3,
    colframe=blue!50!black,
    boxrule=0.8pt,
    arc=2mm
}

\newtcolorbox{important}{
    enhanced,
    breakable,
    title={Importante},
    fonttitle=\bfseries,
    colback=yellow!8,
    colframe=orange!70!black,
    boxrule=0.8pt,
    arc=2mm
}

\newtcolorbox{example}[1]{
    enhanced,
    breakable,
    title={Ejemplo: #1},
    fonttitle=\bfseries,
    colback=green!4,
    colframe=green!40!black,
    boxrule=0.8pt,
    arc=2mm
}

\newtcolorbox{formula}[1]{
    enhanced,
    breakable,
    title={Fórmula / Regla: #1},
    fonttitle=\bfseries,
    colback=purple!4,
    colframe=purple!50!black,
    boxrule=0.8pt,
    arc=2mm
}

\newtcolorbox{exercise}[1]{
    enhanced,
    breakable,
    title={Ejercicio: #1},
    fonttitle=\bfseries,
    colback=gray!5,
    colframe=gray!60!black,
    boxrule=0.8pt,
    arc=2mm
}

\newtcolorbox{note}{
    enhanced,
    breakable,
    title={Nota},
    fonttitle=\bfseries,
    colback=white,
    colframe=gray!50,
    boxrule=0.6pt,
    arc=2mm
}

% ============================================================
% CONFIGURACIÓN FINAL
% ============================================================
\renewcommand{\arraystretch}{1.3}

% ============================================================
\begin{document}
% ============================================================

% ---------------------- PORTADA -----------------------------
\begin{titlepage}

\thispagestyle{empty}

\begin{center}

\vspace*{1cm}

{\large \textsc{\universidad}}

\vspace{3cm}

{\Huge \textbf{\materia}}

\vspace{1cm}

{\LARGE \titulo}

\vspace{0.8cm}

\textit{Comprender los conceptos es el primer paso para convertir el estudio en conocimiento.}

\vspace{1.2cm}

\rule{0.70\textwidth}{0.5pt}

\vspace{0.5cm}

{\large Apuntes de estudio}

\vspace{0.5cm}

\rule{0.70\textwidth}{0.5pt}

\vfill

{\large \autor}

\vspace{0.3cm}

{\large \anio}

\end{center}

\vfill

\begin{flushleft}
\textit{Este es un modesto aporte a los estudiantes de ``Procedimiento Civil'' no pretende sustituir el material oficial de la materia.}
\end{flushleft}

\end{titlepage}

% ---------------------- ÍNDICE ------------------------------
\tableofcontents
\clearpage

% ---------------------- CONTENIDO ---------------------------
\chapter[Cédulas de notificación]{Cédulas de notificación: contexto, importancia y concepto}

\section{Contenido práctico y fundamento normativo}

El contenido práctico de la materia comprende las \textbf{notificaciones}, los \textbf{oficios} y las \textbf{formas de presentar demandas}. Todo lo que se enseña en este plano se fundamenta en los códigos de procedimiento civil y civil argentinos: los instrumentos utilizados son el \textbf{Código Civil} y el \textbf{Código Procesal Civil}, en los cuales se encuentra prácticamente todo lo desarrollado.

\begin{note}
Existen acuerdos y normas adicionales más allá de lo que está codificado, pero la base legal y teórica se encuentra en los instrumentos normativos mencionados.
\end{note}

\section{Importancia del procedimiento de notificaciones}

El dominio del procedimiento de notificaciones resulta fundamental en la formación como profesional del derecho. Las cédulas de notificación son el instrumento clave mediante el cual se garantizan el \textbf{debido proceso} y el \textbf{derecho de defensa} de todas las partes en un litigio.

Comprender no solo la teoría sino también el funcionamiento cotidiano de las notificaciones en los juzgados permite:

\begin{itemize}
    \item gestionar eficientemente los casos como abogado litigante;
    \item evitar errores procedimentales que pueden resultar en \textbf{nulidades};
    \item optimizar tiempos y costos en los procesos;
    \item proteger los derechos de los clientes mediante un correcto seguimiento de las notificaciones;
    \item comprender la realidad administrativa de los tribunales.
\end{itemize}

\section{Concepto de cédula de notificación}

\begin{definition}{Cédula de notificación}
Instrumento formal mediante el cual se comunica a una parte en un proceso judicial sobre actos, resoluciones o citaciones del juzgado.
\end{definition}

La cédula de notificación es un elemento fundamental del \textbf{debido proceso} y garantiza el \textbf{derecho de defensa} de todos los involucrados en el litigio.

% ---------------------- CAPÍTULO 2 --------------------------
\chapter[Cédulas no entregadas y reintento]{Cédulas no entregadas y procedimiento de reintento}

\section{Qué sucede cuando una cédula no es entregada}

Cuando una cédula vuelve sin notificar, el sobre con las copias queda en el juzgado. El oficial de justicia, al no lograr entregar la cédula (por ejemplo, por inasistencia del demandado o por domicilio incorrecto), la devuelve al juzgado junto con el sobre que contenía las copias del acto a notificar.

La cédula devuelta se deposita en un lugar específico del juzgado, comúnmente llamado \textbf{canastilla}, donde queda almacenada para su futura consulta.

% TODO: contenido incompleto o ambiguo en las notas originales (las notas varían entre «canastilla» y «canastillo» al nombrar el recipiente)

\section{Búsqueda de la cédula devuelta en el juzgado}

La búsqueda de las cédulas devueltas (``rebotadas'') es un proceso \textbf{manual}. En la canastilla se acumulan todas las cédulas devueltas, de modo que el profesional debe buscar \textbf{una por una} su propia cédula rebotada entre todas las acumuladas.

\section{Procedimiento de reintento de notificación}

Cuando se decide hacer una nueva cédula \textbf{bajo responsabilidad} o \textbf{con habilitación}, el profesional imprime la cédula y concurre al juzgado. Los pasos indicados en las notas son los siguientes:

\begin{enumerate}
    \item Dirigirse a la \textbf{mesa de entrada} del juzgado.
    \item Solicitar la canastilla de cédulas devueltas.
    \item Buscar manualmente la cédula rebotada propia entre todas las cédulas acumuladas.
    \item Verificar el informe previo del oficial de justicia en el sistema PJN (Poder Judicial en Línea).
    \item Recuperar las copias del acto a notificar.
    \item Preparar las nuevas cédulas, ya impresas y firmadas.
\end{enumerate}

% TODO: contenido incompleto o ambiguo en las notas originales (el orden exacto de los pasos y el significado de la sigla PJN figuran tal como aparecen en las notas)

\section{Optimización del procedimiento de reintento}

Una vez recuperadas las copias, \textbf{no es necesario reimprimir}: las nuevas cédulas, que ya están impresas y firmadas, se agrupan, se abrochan juntas y se dejan nuevamente en el juzgado.

% TODO: contenido incompleto o ambiguo en las notas originales (no se precisa con claridad qué documentos se abrochan entre sí)

Esta práctica representa:

\begin{itemize}
    \item ahorro de tiempo administrativo;
    \item reducción de costos de impresión;
    \item agilización del procedimiento de reintento;
    \item uso eficiente de recursos del juzgado.
\end{itemize}

\begin{important}
Los juzgados guardan las cédulas devueltas por períodos limitados (generalmente dos meses). Si se vuelve a buscar una cédula con retraso, es posible que ya haya sido descartada y será necesario reintentar nuevamente desde el principio.
\end{important}

% ---------------------- CORNELL -----------------------------
\chapter[Cuadros Cornell]{Cuadros Cornell de repaso}

\section{Cédulas de notificación: contexto, importancia y concepto}

\begin{table}[H]
\centering
\begin{tabularx}{\textwidth}{
    >{\raggedright\arraybackslash}p{0.26\textwidth}
    >{\raggedright\arraybackslash}X
}
\toprule
\textbf{Preguntas / palabras clave} &
\textbf{Notas / respuestas} \\
\midrule

\textbf{¿En qué consiste el contenido práctico de la materia y en qué se fundamenta?} &
Comprende las notificaciones, los oficios y las formas de presentar demandas. Se fundamenta en los códigos de procedimiento civil y civil argentinos: el Código Civil y el Código Procesal Civil. \\

\textbf{¿Todo lo enseñado está codificado?} &
La base legal y teórica se encuentra en los códigos, pero existen acuerdos y normas adicionales más allá de lo que está codificado. \\

\textbf{¿Por qué es importante dominar el procedimiento de notificaciones?} &
Porque las cédulas garantizan el debido proceso y el derecho de defensa de todas las partes. Su dominio permite gestionar eficientemente los casos, evitar errores procedimentales que pueden resultar en nulidades, optimizar tiempos y costos, proteger los derechos de los clientes mediante un correcto seguimiento y comprender la realidad administrativa de los tribunales. \\

\textbf{¿Qué es una cédula de notificación y cuál es su función?} &
Es el instrumento formal mediante el cual se comunica a una parte en un proceso judicial sobre actos, resoluciones o citaciones del juzgado. Es un elemento fundamental del debido proceso y garantiza el derecho de defensa de todos los involucrados. \\

\midrule

\multicolumn{2}{p{\dimexpr\textwidth-2\tabcolsep\relax}}{
\textbf{Resumen:}
La cédula de notificación es el instrumento formal de comunicación de actos, resoluciones o citaciones del juzgado, y garantiza el debido proceso y el derecho de defensa. Su correcto manejo evita nulidades y permite proteger los derechos de los clientes. La materia práctica se apoya en el Código Civil y el Código Procesal Civil, además de acuerdos y normas adicionales.
} \\

\bottomrule
\end{tabularx}
\end{table}

\section{Cédulas no entregadas y procedimiento de reintento}

\begin{table}[H]
\centering
\begin{tabularx}{\textwidth}{
    >{\raggedright\arraybackslash}p{0.26\textwidth}
    >{\raggedright\arraybackslash}X
}
\toprule
\textbf{Preguntas / palabras clave} &
\textbf{Notas / respuestas} \\
\midrule

\textbf{¿Qué sucede cuando la cédula no puede entregarse?} &
El oficial de justicia la devuelve al juzgado junto con el sobre que contenía las copias del acto a notificar (por ejemplo, por inasistencia del demandado o domicilio incorrecto). La cédula devuelta se deposita en la canastilla del juzgado. \\

\textbf{¿Cómo se busca una cédula rebotada?} &
Es un proceso manual: el profesional solicita la canastilla en la mesa de entrada y busca una por una su cédula entre todas las cédulas devueltas acumuladas. \\

\textbf{¿Cuál es el procedimiento para reintentar una notificación?} &
(1) Ir a la mesa de entrada del juzgado; (2) solicitar la canastilla de cédulas devueltas; (3) buscar manualmente la cédula propia; (4) verificar el informe previo del oficial de justicia en el sistema PJN; (5) recuperar las copias del acto a notificar; (6) preparar las nuevas cédulas, ya impresas y firmadas. \\

\textbf{¿Cómo se optimiza el reintento?} &
Sin reimprimir: las nuevas cédulas, ya impresas y firmadas, se agrupan, se abrochan y se dejan nuevamente en el juzgado. Esto implica ahorro de tiempo administrativo, reducción de costos de impresión, agilización del procedimiento y uso eficiente de recursos del juzgado. \\

\textbf{¿Qué riesgo existe si la búsqueda se retrasa?} &
Los juzgados guardan las cédulas devueltas por períodos limitados (generalmente dos meses); si la búsqueda se demora, la cédula puede haber sido descartada y habrá que reintentar desde el principio. \\

\midrule

\multicolumn{2}{p{\dimexpr\textwidth-2\tabcolsep\relax}}{
\textbf{Resumen:}
Cuando una cédula no se entrega, vuelve al juzgado con las copias y queda en la canastilla. Para reintentar, el profesional busca manualmente su cédula, verifica el informe del oficial de justicia en el PJN, recupera las copias y prepara nuevas cédulas ya impresas y firmadas, que se agrupan y se dejan sin reimprimir. Debe actuar con rapidez, pues las cédulas devueltas se conservan generalmente dos meses.
} \\

\bottomrule
\end{tabularx}
\end{table}

\end{document}
